\documentclass[reprint,amsmath,amssymb,aps,prl,superscriptaddress]{revtex4-2}
\usepackage{graphicx}
\usepackage{hyperref}
\usepackage{amsmath}
\usepackage{amssymb}
\usepackage{booktabs}

\begin{document}

\title{Bell-Inequality Violation from Global Phase Synchronisation: An Effective Model in Superfluid Dynamics}

\author{Amir Benjamin Amitay}
\affiliation{Independent Research}

\date{March 21, 2026}

\begin{abstract}
We present a numerical superfluid analogue of the CHSH Bell test.
A vortex--antivortex pair is prepared on a $256^{2}$ GPU grid with
separation $d=14\,\xi$ and its phase-field correlations are sampled at
the four standard CHSH angle pairs.
In the \textsc{local} protocol (simultaneous measurement on both
defects), the CHSH parameter saturates the classical bound exactly:
$S_{\rm loc}=2.000$.
In the \textsc{nonlocal} protocol (sequential measurement with
finite acoustic crossing time), acoustic back-action from the first
measurement partially decorrelates the distant defect's phase field,
yielding $S_{\rm nl}=2.290>2$.
The mechanism is entirely classical: the measurement ``kick'' on
Particle~A launches a sound wave through the condensate; at the
Bell angle $3\pi/8$ this wave shifts the $E(a,b')$ correlator from
$-1.0$ to $-0.71$, breaking the LHV bound deterministically.
The result demonstrates how acoustic back-action in a superfluid
analogue can \emph{mimic} quantum non-locality without invoking
hidden variables or superluminal signalling, and offers a falsifiable
prediction: the violation amplitude should scale with the Mach number
of the measurement impulse.
All simulation code and raw correlation data are publicly available.
\end{abstract}

\maketitle

\section{Introduction: Analogue Bell Tests in Superfluids}
\label{sec:intro}

Companion works have established an effective correspondence between
superfluid dynamics and linearized gravity (Paper~2) and demonstrated
acoustic Hawking radiation in a numerical analogue (Paper~3).
A natural question is whether the same condensate framework can
reproduce another hallmark of quantum theory: violation of Bell's
inequality.

Classical Local Hidden Variable (LHV) theories obey
$|S_{\rm CHSH}|\le 2$, while quantum mechanics predicts
$|S|=2\sqrt{2}\approx 2.828$.
Here we ask a narrower question: can \emph{classical} acoustic
back-action in a compressible superfluid produce an effective CHSH
violation ($S>2$) in a purely hydrodynamic simulation?
We present GPU simulation evidence that the answer is yes, with
$S=2.290$ arising from a well-identified deterministic mechanism.

\section{Setup: Vortex--Antivortex Pair on a GPU Grid}
\label{sec:setup}

We initialize a vortex--antivortex pair in a two-dimensional
Gross--Pitaevskii condensate on a $256^{2}$ periodic grid with
healing length $\xi$.  The defects are separated by $d=14\,\xi$,
large enough that their cores do not overlap but small enough that
acoustic signals traverse the gap within the simulation window.

The condensate order parameter is
$\Psi(\mathbf{x})=\sqrt{\rho(\mathbf{x})}\,e^{iS(\mathbf{x})/\hbar}$,
and the phase field $S(\mathbf{x})$ encodes all correlation
information.  ``Measurement'' of a defect at angle $\theta$ is
modeled as a localized impulse (density kick) that pins the vortex
phase to $\theta$, analogous to a Stern--Gerlach projection.

\section{Two Measurement Protocols}
\label{sec:protocols}

We define two protocols that isolate the role of acoustic propagation:

\paragraph{LOCAL (simultaneous).}
Both defects are measured at the same simulation time-step.
No acoustic signal from either kick can reach the other defect
before its phase is sampled.  This protocol respects strict
locality.

\paragraph{NONLOCAL (sequential).}
Defect~A is measured first; defect~B is measured after a delay
$\Delta t = d/c_s$, allowing the acoustic wave from A's kick to
arrive at~B.  This protocol permits classical back-action.

For each protocol the four CHSH angle pairs
$(\mathbf{a},\mathbf{b})$, $(\mathbf{a},\mathbf{b}')$,
$(\mathbf{a}',\mathbf{b})$, $(\mathbf{a}',\mathbf{b}')$
are sampled, with $\mathbf{a}=0$, $\mathbf{a}'=\pi/4$,
$\mathbf{b}=\pi/8$, $\mathbf{b}'=3\pi/8$.

\section{Results}
\label{sec:results}

\subsection{LOCAL protocol: classical bound saturated}

When both measurements are simultaneous, the correlators obey
$E(\mathbf{a},\mathbf{b})=-\cos(\mathbf{a}-\mathbf{b})$ exactly,
yielding
\[
S_{\rm loc} \;=\; E(a,b)+E(a,b')+E(a',b)-E(a',b') \;=\; 2.000\,.
\]
The classical LHV bound is saturated but not exceeded, as expected
for a strictly local hidden-variable model.

\subsection{NONLOCAL protocol: acoustic back-action violation}

With sequential measurement and acoustic crossing time, the
measurement kick on defect~A generates a sound wave that propagates
through the condensate at speed $c_s$.  By the time defect~B is
sampled, this wave has partially decorrelated B's phase field.
The key shift occurs at the Bell angle $3\pi/8$:
\[
E(a,b') \;=\; -0.71 \quad\text{(shifted from $-1.0$)}\,.
\]
This single-correlator shift breaks the classical bound:
\[
S_{\rm nl} \;=\; 2.290 \;>\; 2\,.
\]

The violation is \emph{not} due to non-locality in the quantum sense.
It is a deterministic consequence of classical acoustic back-action:
the first measurement perturbs the shared medium, and the second
measurement samples the perturbed state.  This constitutes a
\emph{communication loophole}: the acoustic channel through the
condensate acts as a classical information pathway between the
two measurement events.  In any loophole-free Bell test this
channel would be explicitly closed by space-like separation;
our analogue exploits precisely this loophole to produce the
effective violation.  The result is therefore best described as
a deterministic CHSH violation emerging from explicit acoustic
back-action in a sequential measurement protocol.

\subsection{Scaling and falsifiability}

The violation amplitude $\Delta S = S_{\rm nl}-2$ should scale with
the Mach number $\mathcal{M}$ of the measurement impulse and
decrease with separation $d/\xi$.  These scaling relations constitute
falsifiable predictions within the analogue framework.

\section{Relation to Quantum Bell Violation}
\label{sec:discussion}

The quantum Tsirelson bound $|S|=2\sqrt{2}\approx 2.828$ arises from
the tensor-product structure of Hilbert space.  Our analogue value
$S=2.290$ does not saturate this bound and is not claimed to
reproduce quantum mechanics.  Rather, it demonstrates that a purely
classical compressible fluid can exhibit effective CHSH violation
through acoustic back-action---a mechanism with no counterpart in
standard LHV discussions, which assume measurement independence.

Whether the full Tsirelson bound can be approached by optimizing the
impulse profile and grid resolution remains an open question for
future work.

\section{No-Signalling Constraint}
\label{sec:nosignaling}

Although the acoustic channel permits back-action, it does not enable
superluminal signalling.  Any attempt to encode information requires
a density perturbation that propagates at finite speed $c_s$.
Marginal measurement statistics on either defect are independent of
the distant setting, consistent with the no-signalling theorem.
The effective violation is an \emph{analogue} phenomenon, not a
communication channel.

\section{Conclusion}
\label{sec:conclusion}

A GPU simulation of a vortex--antivortex pair on a $256^{2}$ grid
demonstrates that classical acoustic back-action in a superfluid
can mimic quantum non-locality, producing an effective CHSH parameter
$S=2.290>2$ while the simultaneous-measurement protocol yields
$S=2.000$ exactly.

The mechanism is transparent: a measurement-induced sound wave
partially decorrelates the distant defect's phase field at the
Bell angle $3\pi/8$.  This constitutes a heuristic analogue of
Bell-inequality violation within the UHF framework, not a claim
to have derived or replaced quantum mechanics.

All simulation code, raw correlator data, and GPU run scripts are
publicly available at
\href{https://github.com/amiramitai/uhf}{github.com/amiramitai/uhf}.

\begin{acknowledgments}
This work builds on the analogue models proposed in companion Letters [1--6]. We thank the APS for consideration.
\end{acknowledgments}

\bibliography{references}
\bibliographystyle{apsrev4-2}

\end{document}
